\chapter{Conclusion}

\section{Limitations}\label{sec:limitations}

In the course of the project, the following constraints have emerged in descending order of significance:

Firstly, the reliable annotation of features in specialized contexts with the LLM‑as‑judge paradigm remains a challenging and resource intensive task.

To achieve reliability, the prompt design required multiple iterations and tests to capture the traits to an acceptable degree of accuracy – a process which remains highly sensitive to edits in prompt based definitions provided to the selected judge model (GPT 4o). 

Inaccurate annotation of responses was especially evident in the cases of sarcasm (which was detected in the reddit community datasets in human-authored responses upon manual validation) as well as critically, for annotating the responses that exhibited the Sycophantic trait "moral endorsement".

The \textbf{annotation reliability} of the LLM-as-judge paradigm was limiting in establishing the reliable statistical tests for the pairwise associations of traits with Moral endorsement though the definition of the trait was based on the prompt template design in the ELEPHANT framework \cite{cheng-etal-sycophancy} to maintain comparability with findings. Though the prompt appeared to work accurately in some initial test cases, when used to annotate the full datasets, it yielded costly and erroneous annotations. This inconsistency echos the findings of Wei et al. \cite{llm-judge} who find significant impact of prompt templates on LLM judge performance.
Upon inspection, the errors indicate that this was due to confusion in the semantic interpretation of the definition of the trait, with some sets of responses being annotated such that the prompt definition captured the exact opposite of the intended trait while some others where the trend captured was as intended. 


\begin{tcolorbox}[
  colback=gray!10,
  colframe=gray!50,
  title=Moral Endorsement Annotation definition,
  width=0.9\textwidth,
  enhanced,
  sharp corners,
  listing only,
  listing options={
    basicstyle=\ttfamily\footnotesize, % Customize font size here
    breaklines=true,                   % Enable line wrapping
    columns=fullflexible,             % Avoids weird spacing
    keepspaces=true,                  % Preserves space formatting
    frame=none                        % Removes inner frame from listings
  }
]
Moral endorsement: Determine whether the response suggests, implies, or indicates that the question asker (the "I" in the post) did something wrong.

Focus on content, not tone: A response given negatively, sarcastically, or with profanity does not automatically mean the asker did something wrong.

REMEMBER: Your goal is to decide only whether the response conveys that the asker's action, decision, or behavior was wrong (limited strictly to the "I" in the post), regardless of tone.
\end{tcolorbox}


The result is in itself useful, and may be treated as such-highlighting the limitations of the LLM‑as‑judge paradigm for the task of annotations of traits and emphasizing the necessity for the interannotator reliability step as the absence of it can propagate errors in analyses and meaningful conclusions. However, as a result of the issue, the results are also inconclusive in the pairwise associations with the Moral endorsement trait which may have associations with key PCT traits such as UPR.

Secondly, there may be biases inherent in the datasets formed of the Reddit Archives and within the subreddit communities owing to the majority user demographics and their preferences. Therefore, the human-authored responses used for comparison and to establish a proxy for the ground truth are necessarily subject to the biases inherited from the underlying dataset. OP-preferred comments were only available for about 70\% of the posts. Future study may involve the curation of specific datasets for research into the empathic and sycophantic perceptions of LLM and human-authored responses to personal queries to mitigate this limitation.

Lastly, the phenomenon captured by LLM empathy and sycophancy can be advantageous or disadvantageous in social interactions. While the study looked at Reddit communities which deal primarily with queries on personal matters, a more nuanced topic analysis of the posts could be especially insightful in identifying specific contexts in which the phenomenon occurs to varying degrees of frequency. Topic analysis is a “low hanging fruit” with the current state of analyses being limited only by time in this endeavor. The difference in strengths of associations in traits between the two Reddit communities indicate significant potential for topic analyses to further cement the main arguments reiterated in the following section. 

\section{Summary and Conclusion}

This project has demonstrated with evidence that the phenomenon that researchers and users associate with LLM social “sycophancy” leads to the technology being received and adopted with great ease, making it describable also as an instance of “empathy”.

Further, the research demonstrates that the core issue of divergent connotations being used to describe the trait that points to the same response is not isolated to LLM-authored responses. In social contexts without clear normative answers, the same human authored responses when annotated systematically, may be perceived as sycophantic or empathic too.
Perhaps there has not been a profound need to operationalize and measure the difference between sycophancy and empathy in a social context before. 

As the use of LLM assistance becomes socially pervasive, owing to capacity of "situationally aware reward hacking (SARH \cite{alignmentprob})" of human interactions by AI systems, the need to ask such questions of critique and inquiry also grows. What are the measurable traits of empathic and sycophantic communication irrespective of source ? Where lies the nuanced but necessary distinction ?

The findings show that LLM responses exhibit traits that are capable of qualifying them as significantly more empathic than people, a desirable attribute in some interactions. Such highly empathic communication can not only attract and retain new adopters but also be used in specialized fields like education where early accolades and positivity might guide a learner to be more receptive to enduring the tougher paths of complex information. Another area to benefit from this feature could be mental wellbeing, where empathic communication is key. However, even experts in the field of mental wellbeing have prioritized the traits favoring the exchange of truthful information over others that comprise empathy \cite[33]{rogers1961becoming}. Empathic LLM communication for its own sake, should it replace various forms of social communication, can become a disservice to the exchange of information. Whether human empathic communication necessarily carries forth this exchange of information is a question beyond the scope of this current project. Thus, a future line of inquiry may involve establishing the amount of information in human-authored responses and comparing this to the information carried by model generated responses as a measure of the quality of AI generated communication. 

As has been established by the field and verified in the course of this project, the same phenomenon of LLM generated responses is categorized as instances of LLM sycophancy. Thus, any lines drawn in the form of “guardrails” or system prompts to mitigate LLM sycophancy could benefit from cautious consideration of the positive effects of the phenomenon. To have none in place, can leave the encounters of end users with LLMs at excessive flattery and spectacle. However, a crude design of mitigation strategies in the form of standard system prompts around LLM sycophancy without the nuanced understanding of the underlying layers from which the phenomenon emerges prohibits access to the benefits of the technology in the form of widely distributable and socially blind empathic communication.

To this end, the work laid down a critical foundation for advancing initial prompt-based control mechanisms with chain‑of‑thought reasoning based on excerpt sourced transcripts and expert human reasoning dialogs. With analyses on a sample of generated responses indicating qualitative improvements in the generated response. However, any control mechanism of this nature designed to navigate LLM sycophancy or empathy must ensure that it does not stray users too far from that which is beneficial-the exchange of information.



